% !TEX root = ../main.tex

\chapter{Mission principale~: l'historisation des modifications}
\label{chap:historisation}

\section{Contexte et besoin métier}

Le FPG traite quotidiennement un grand nombre de dossiers de financement,
impliquant des montants financiers conséquents et des données sensibles
concernant les structures adhérentes et leurs collaborateurs. Dans ce
contexte, la capacité à \textbf{tracer et afficher l'historique complet
des modifications} apportées à un dossier constitue un besoin métier
essentiel.

Concrètement, il s'agissait de permettre, pour chaque demande de
financement, chaque structure adhérente et chaque compte collaborateur,
de répondre à trois questions pour chaque modification apportée :
\textbf{qui} a effectué la modification, \textbf{quand}, et
\textbf{quelle était l'ancienne valeur par rapport à la nouvelle}. Cette
fonctionnalité répond à des enjeux de :

\begin{itemize}
  \item \textbf{traçabilité} — savoir précisément ce qui a été modifié
        et par qui à chaque étape du traitement d'un dossier ;
  \item \textbf{audit} — permettre un contrôle a posteriori du
        traitement des demandes de financement ;
  \item \textbf{résolution de litiges} — disposer de preuves fiables en
        cas de contestation d'une décision ou d'une donnée par une
        structure adhérente.
\end{itemize}

\section{Choix techniques}

\subsection{django-simple-history}

Le choix s'est porté sur la bibliothèque \texttt{django-simple-history},
qui génère automatiquement une table \texttt{\_history} pour chaque
modèle historisé. Chaque opération \texttt{INSERT}, \texttt{UPDATE} ou
\texttt{DELETE} effectuée sur un modèle historisé donne lieu à
l'enregistrement d'une nouvelle ligne dans cette table, capturant
l'intégralité de l'état de l'objet à cet instant, ainsi que l'utilisateur
responsable de la modification et l'horodatage associé.

\subsection{Le \texttt{HistoryService}}

Afin d'éviter la duplication de logique entre les différentes entités
historisées, j'ai construit un service dédié, le \texttt{HistoryService},
centralisant la récupération et la mise en forme de l'historique pour
chaque type d'entité. Ce service constitue le point d'entrée unique
utilisé par les différents endpoints d'API exposant de l'historique.

\begin{lstlisting}[language=Python, caption={Extrait simplifié du \texttt{HistoryService}}]
class HistoryService:
    def get_history(self, entity, entity_id):
        history_records = self._fetch_history(entity, entity_id)
        resolved = self._resolve_references(history_records)
        return self._format_for_frontend(resolved)
\end{lstlisting}

\subsection{Système de snapshots}

Une demande de financement regroupe de nombreuses entités liées
(stagiaires, allocations budgétaires, tiers bénéficiaires, règlements).
Comparer indépendamment l'historique de chacune de ces entités aurait
rendu difficile la reconstitution d'un état cohérent du dossier à un
instant donné.

J'ai donc mis en place un système de \textbf{snapshots}, matérialisé par
le modèle \texttt{FundingRequestSnapshot}, capturant à chaque sauvegarde
un état complet et cohérent de l'ensemble des entités liées à une demande
de financement. La comparaison entre deux snapshots successifs, via une
méthode \texttt{diff\_against}, permet d'obtenir de façon fiable
l'ensemble des changements survenus entre deux états du dossier.

\section{Logique métier et historisation par module fonctionnel}
\label{sec:historisation-par-onglet}

La demande de financement se compose de plusieurs onglets, chacun
correspondant à un domaine fonctionnel distinct et porteur d'une logique
d'historisation propre. Cette section détaille, pour chaque onglet, les
informations concernées et la manière dont leurs modifications sont
tracées et restituées à l'utilisateur.

\subsection{Onglet Formation}

\begin{figure}[H]
  \centering
  \imgw{onglet_formation.png}{0.9\textwidth}
  \caption{Onglet Formation}
  \label{fig:onglet-formation}
\end{figure}

L'onglet Formation regroupe les informations générales relatives à
l'action de formation~: le type de prestation (action de formation,
bilan de compétences, validation des acquis de l'expérience, permis de
conduire), la localisation en Polynésie ou non, le caractère externe ou
interne du prestataire, l'enregistrement au Service de l'Emploi, de la
Formation et de l'Insertion Professionnelles (SEFI), ainsi que le nom du
formateur. S'y ajoutent les informations propres à la formation
elle-même~: son nom, sa modalité (présentiel ou distanciel), ses
domaines et sous-domaines, et son lieu de déroulement (chez le
prestataire, dans les locaux de l'entreprise, ou ailleurs).

Certains champs de cet onglet présentent une particularité importante~:
leur modification déclenche un recalcul automatique du montage financier
du dossier, un mécanisme désigné en interne sous le terme de
\textbf{ventilation}. Il s'agit notamment~:

\begin{itemize}
  \item de la sélection des dates et heures de formation~;
  \item du coût pédagogique, du nom du cofinanceur et du montant du
        cofinancement~;
  \item des dispositifs complémentaires associés au dossier.
\end{itemize}

\begin{figure}[H]
  \centering
  \imgw{ajout_dispositif_complementaire.png}{0.8\textwidth}
  \caption{Ajout d'un dispositif complémentaire}
  \label{fig:ajout-dispositif}
\end{figure}

Cette distinction entre champs \og{}simples\fg{} et champs déclencheurs
de ventilation a directement structuré la conception du système
d'historisation~: pour ces derniers, il ne suffisait pas de tracer le
changement de valeur sur l'onglet Formation, mais également son
répercussion sur le montage financier (voir
Section~\ref{subsec:montage-financier}).

\subsection{Onglet Stagiaires}

\begin{figure}[H]
  \centering
  \imgw{onglet_stagiaire.png}{0.9\textwidth}
  \caption{Onglet Stagiaires}
  \label{fig:onglet-stagiaires}
\end{figure}

Les stagiaires correspondent aux salariés de la structure adhérente
participant à la formation. Plusieurs types de modifications devaient
être historisés sur cet onglet.

\paragraph{Ajout et suppression d'un stagiaire.}
L'ajout comme la suppression d'un stagiaire du dossier constituent des
événements à part entière, affichés distinctement dans l'historique.

\begin{figure}[H]
  \centering
  \imgw{Ajout_stagiaire.png}{0.75\textwidth}
  \caption{Ajout d'un stagiaire, vue formulaire}
  \label{fig:ajout-stagiaire}
\end{figure}

\begin{figure}[H]
  \centering
  \imgw{ajout_stagiaire_historique.png}{0.75\textwidth}
  \caption{Ajout d'un stagiaire, restitution dans l'historique}
  \label{fig:ajout-stagiaire-hist}
\end{figure}

\begin{figure}[H]
  \centering
  \imgw{suppression_stagiaire.png}{0.75\textwidth}
  \caption{Suppression d'un stagiaire}
  \label{fig:suppression-stagiaire}
\end{figure}

\paragraph{Modification des informations d'un stagiaire.}
La modification du type de contrat, des dates de début et de fin de
contrat, du niveau de formation, de la fonction occupée ou de la
catégorie socio-professionnelle (CSP) du stagiaire fait également l'objet
d'un suivi.

\begin{figure}[H]
  \centering
  \imgw{modification_stagiaire.png}{0.75\textwidth}
  \caption{Modification des informations d'un stagiaire}
  \label{fig:modif-stagiaire}
\end{figure}

\begin{figure}[H]
  \centering
  \imgw{modification_stagiaire_historique.png}{0.75\textwidth}
  \caption{Restitution dans l'historique}
  \label{fig:modif-stagiaire-hist}
\end{figure}

\paragraph{Modification du salaire et des heures travaillées.}
Le salaire brut mensuel et le nombre d'heures travaillées par mois
constituent des données particulières~: comme les champs évoqués
précédemment sur l'onglet Formation, leur modification a un impact direct
sur le calcul du montage financier du stagiaire concerné.

\begin{figure}[H]
  \centering
  \imgw{modification_stagiaire_salaire_heures.png}{0.75\textwidth}
  \caption{Modification du salaire et des heures travaillées}
  \label{fig:modif-salaire-heures}
\end{figure}

\paragraph{Répercussion sur le montage financier.}
Pour ces trois catégories d'action (ajout, suppression, modification du
salaire ou des heures), il était donc nécessaire d'afficher, en plus de
l'événement lui-même, le détail de la ventilation qu'il entraîne sur le
montage financier du dossier.

\begin{figure}[H]
  \centering
  \imgw{ajout_stagiaire_ventilation.png}{0.75\textwidth}
  \caption{Ventilation entraînée par l'ajout d'un stagiaire}
  \label{fig:vent-ajout-stagiaire}
\end{figure}

\begin{figure}[H]
  \centering
  \imgw{supression_stagiaire_ventilation.png}{0.75\textwidth}
  \caption{Ventilation entraînée par la suppression d'un stagiaire}
  \label{fig:vent-suppr-stagiaire}
\end{figure}

\begin{figure}[H]
  \centering
  \imgw{modification_stagiaire_salaire_ventilation.png}{0.75\textwidth}
  \caption{Ventilation entraînée par la modification du salaire}
  \label{fig:vent-salaire}
\end{figure}

\subsection{Onglet Suivi de formation}

\begin{figure}[H]
  \centering
  \imgw{onglet_suivie_de_formation.png}{0.9\textwidth}
  \caption{Onglet Suivi de formation}
  \label{fig:onglet-suivi}
\end{figure}

Cet onglet répertorie les heures de formation effectivement réalisées
par chaque stagiaire. Il présente la particularité de n'être modifiable
que lorsque le dossier se trouve au statut \og{}Préparation
règlement\fg{}, conformément à la matrice de verrouillage présentée au
chapitre précédent. Comme pour les champs évoqués plus haut, la
modification des heures réalisées influence directement le montage
financier du dossier et déclenche par conséquent une ventilation, elle
aussi restituée dans l'historique.

\begin{figure}[H]
  \centering
  \imgw{modification_suivie_de_formation.png}{0.75\textwidth}
  \caption{Modification des heures réalisées}
  \label{fig:modif-suivi}
\end{figure}

\begin{figure}[H]
  \centering
  \imgw{modification_suivie_de_formation_ventilation.png}{0.75\textwidth}
  \caption{Ventilation entraînée par la modification des heures réalisées}
  \label{fig:vent-suivi}
\end{figure}

\subsection{Onglet Montage financier}
\label{subsec:montage-financier}

\begin{figure}[H]
  \centering
  \imgw{onglet_montage_financier.png}{0.9\textwidth}
  \caption{Onglet Montage financier}
  \label{fig:onglet-montage}
\end{figure}

L'onglet Montage financier répertorie, pour chaque stagiaire et chaque
type de frais, la source de financement retenue~: selon les dispositifs
complémentaires activés sur le dossier, certains frais peuvent être pris
en charge par différentes lignes budgétaires.

Deux modes de modification du montage financier devaient être distingués
dans l'historique~:

\begin{itemize}
  \item une modification \textbf{par ventilation}, c'est-à-dire une
        modification \textit{indirecte}, déclenchée par un changement
        apporté sur un autre onglet (dates de formation, coût
        pédagogique, dispositifs complémentaires, stagiaires, heures
        réalisées)~;
  \item une modification \textbf{manuelle}, apportée directement sur
        l'onglet du montage financier par un collaborateur.
\end{itemize}

\begin{figure}[H]
  \centering
  \imgw{modification_montage_financer_manuel.png}{0.75\textwidth}
  \caption{Modification manuelle du montage financier}
  \label{fig:montage-manuel}
\end{figure}

\begin{figure}[H]
  \centering
  \imgw{modification_montage_financer_manuel_historique.png}{0.75\textwidth}
  \caption{Restitution d'une modification manuelle dans l'historique}
  \label{fig:montage-manuel-hist}
\end{figure}

Le volume d'informations généré par une ventilation — potentiellement
plusieurs lignes de dispositif, pour plusieurs stagiaires, sur plusieurs
types de frais — imposait un traitement différencié de son affichage.
Deux exigences ont guidé la conception de cette restitution~:

\begin{enumerate}
  \item afficher un \textbf{événement source}, c'est-à-dire indiquer
        clairement quelle modification, sur quel onglet, est à l'origine
        de la ventilation observée~;
  \item regrouper le détail de la ventilation \textbf{par stagiaire},
        avec un mécanisme de dépliage à la demande, afin de ne pas
        surcharger l'affichage par défaut.
\end{enumerate}

\begin{figure}[H]
  \centering
  \imgw{modification_montage_financer_ventilation_historique.png}{0.8\textwidth}
  \caption{Restitution d'une ventilation dans l'historique, avec
    événement source et regroupement par stagiaire}
  \label{fig:montage-ventilation-hist}
\end{figure}

\subsection{Onglet Tiers bénéficiaires}

\begin{figure}[H]
  \centering
  \imgw{onglet_tiers_beneficiaires.png}{0.9\textwidth}
  \caption{Onglet Tiers bénéficiaires}
  \label{fig:onglet-tiers}
\end{figure}

L'onglet Tiers bénéficiaires permet de déterminer, pour chaque type de
frais, lequel de l'adhérent, de l'organisme de formation ou du FPG en est
le bénéficiaire. L'historisation de cet onglet, plus simple que celle du
montage financier, consiste à tracer la modification du bénéficiaire
retenu pour chaque catégorie de frais.

\begin{figure}[H]
  \centering
  \imgw{modification_tiers_bénéficiaire.png}{0.75\textwidth}
  \caption{Modification d'un tiers bénéficiaire}
  \label{fig:modif-tiers}
\end{figure}

\begin{figure}[H]
  \centering
  \imgw{modification_tiers_bénéficiaire_historique.png}{0.75\textwidth}
  \caption{Restitution dans l'historique}
  \label{fig:modif-tiers-hist}
\end{figure}

\subsection{Onglet Règlements}

\begin{figure}[H]
  \centering
  \imgw{onglet_règlements.png}{0.9\textwidth}
  \caption{Onglet Règlements}
  \label{fig:onglet-reglements}
\end{figure}

L'onglet Règlements n'est modifiable que lorsque le dossier se trouve
dans l'un des statuts \og{}Préparation règlement\fg{}, \og{}En attente
conformité règlement\fg{} ou \og{}Réglé\fg{}. Trois types d'événements
devaient y être historisés~:

\begin{itemize}
  \item l'\textbf{ajout} d'un règlement~;
  \item la \textbf{suppression} d'un règlement~;
  \item la \textbf{validation} ou l'\textbf{invalidation} d'un
        règlement.
\end{itemize}

Une règle métier importante encadre ce dernier point~: un règlement déjà
validé ne peut pas être supprimé directement — il doit au préalable être
invalidé. Cette contrainte fonctionnelle a dû être prise en compte lors
des campagnes de tests de l'historisation, afin de vérifier que chaque
transition d'état d'un règlement est correctement et distinctement
tracée.

\begin{figure}[H]
  \centering
  \imgw{ajout_règlement.png}{0.75\textwidth}
  \caption{Ajout d'un règlement}
  \label{fig:ajout-reglement}
\end{figure}

\begin{figure}[H]
  \centering
  \imgw{ajout_règlement_historique.png}{0.75\textwidth}
  \caption{Restitution de l'ajout dans l'historique}
  \label{fig:ajout-reglement-hist}
\end{figure}

\begin{figure}[H]
  \centering
  \imgw{suppression_règlement.png}{0.75\textwidth}
  \caption{Suppression d'un règlement}
  \label{fig:suppr-reglement}
\end{figure}

\begin{figure}[H]
  \centering
  \imgw{supression_règlement_historique.png}{0.75\textwidth}
  \caption{Restitution de la suppression dans l'historique}
  \label{fig:suppr-reglement-hist}
\end{figure}

\begin{figure}[H]
  \centering
  \imgw{validation_invalidation_règlement_historique.png}{0.75\textwidth}
  \caption{Restitution de la validation/invalidation d'un règlement}
  \label{fig:valid-reglement-hist}
\end{figure}

\subsection{Onglet Documents}

\begin{figure}[H]
  \centering
  \imgw{onglet_document.png}{0.9\textwidth}
  \caption{Onglet Documents}
  \label{fig:onglet-documents}
\end{figure}

Contrairement aux onglets précédents, l'onglet Documents reste
modifiable quel que soit le statut du dossier. Il distingue deux
catégories de fichiers, dont la nature différente a directement
conditionné la logique d'historisation retenue.

\paragraph{Les rapports (\textit{Reports}).}
Il s'agit des documents générés automatiquement par le système, en
particulier les \textit{MAPC}, produits lorsque le dossier atteint le
statut \og{}Préparation règlement\fg{}. La génération de ce document est
notamment un préalable obligatoire à toute modification des règlements.
S'agissant de documents officiels générés par le système, seuls leur
ajout et leur visibilité sont historisés — ils ne peuvent en particulier
faire l'objet ni d'une suppression, ni d'une labellisation.

\begin{figure}[H]
  \centering
  \imgw{ajout_report_history.png}{0.75\textwidth}
  \caption{Ajout d'un rapport (MAPC), restitution dans l'historique}
  \label{fig:ajout-report}
\end{figure}

\begin{figure}[H]
  \centering
  \imgw{visibilite_report.png}{0.75\textwidth}
  \caption{Modification de la visibilité d'un rapport}
  \label{fig:visibilite-report}
\end{figure}

\begin{figure}[H]
  \centering
  \imgw{visibilié_report_historique.png}{0.75\textwidth}
  \caption{Restitution de la visibilité d'un rapport dans l'historique}
  \label{fig:visibilite-report-hist}
\end{figure}

\paragraph{Les documents.}
Il s'agit des pièces justificatives importées manuellement par
l'utilisateur, relatives à la formation — par exemple un justificatif de
frais d'hébergement. Pour cette catégorie, l'ensemble du cycle de vie du
document devait être historisé~: son ajout, sa suppression, la
modification de sa visibilité, son verrouillage ou déverrouillage, ainsi
que l'ajout ou le retrait de labels permettant de le catégoriser (par
exemple le label \og{}contrat de formation\fg{}).

\begin{figure}[H]
  \centering
  \imgw{ajout_document_historique.png}{0.75\textwidth}
  \caption{Ajout d'un document, restitution dans l'historique}
  \label{fig:ajout-document-hist}
\end{figure}

\begin{figure}[H]
  \centering
  \imgw{supression_document_historique.png}{0.75\textwidth}
  \caption{Suppression d'un document, restitution dans l'historique}
  \label{fig:suppr-document-hist}
\end{figure}

\begin{figure}[H]
  \centering
  \imgw{visibilité_document_historique.png}{0.75\textwidth}
  \caption{Modification de la visibilité d'un document}
  \label{fig:visibilite-document-hist}
\end{figure}

\begin{figure}[H]
  \centering
  \imgw{verouillage_deverouillage_document_historique.png}{0.75\textwidth}
  \caption{Verrouillage / déverrouillage d'un document}
  \label{fig:verrou-document-hist}
\end{figure}

\begin{figure}[H]
  \centering
  \imgw{ajout_label_document.png}{0.75\textwidth}
  \caption{Ajout d'un label sur un document}
  \label{fig:ajout-label}
\end{figure}

\begin{figure}[H]
  \centering
  \imgw{ajout_label_document_historique.png}{0.75\textwidth}
  \caption{Restitution de l'ajout d'un label dans l'historique}
  \label{fig:ajout-label-hist}
\end{figure}

\begin{figure}[H]
  \centering
  \imgw{supression_label_historique.png}{0.75\textwidth}
  \caption{Restitution de la suppression d'un label dans l'historique}
  \label{fig:suppr-label-hist}
\end{figure}

\subsection{Onglet Historique}

\begin{figure}[H]
  \centering
  \imgw{onglet_historique.png}{0.9\textwidth}
  \caption{Onglet Historique}
  \label{fig:onglet-historique}
\end{figure}

L'onglet Historique constitue l'aboutissement fonctionnel de la
mission~: il centralise et restitue l'ensemble des modifications
apportées à la demande de financement, qu'elles soient issues d'une
action manuelle ou d'une ventilation, tous onglets confondus.

Chaque entrée de l'historique respecte un même format de présentation,
garantissant une lecture homogène quel que soit le type de modification
concernée~:

\begin{quote}
\textit{Modification -- \og{}onglet concerné\fg{}}\\
\textit{Mise à jour de \og{}nom de l'utilisateur\fg{} le \og{}date de
la modification\fg{}}\\
\textit{\og{}type de modification\fg{} de \og{}ancienne valeur\fg{} vers
\og{}nouvelle valeur\fg{}}
\end{quote}

Au-delà des modifications de champs classiques, l'onglet Historique
restitue également deux événements particuliers~: le changement de
statut de la demande de financement, et la création initiale du dossier.

\begin{figure}[H]
  \centering
  \imgw{modification_etat_historique.png}{0.75\textwidth}
  \caption{Restitution d'un changement de statut dans l'historique}
  \label{fig:modif-etat-hist}
\end{figure}

\begin{figure}[H]
  \centering
  \imgw{creation_demande_de_financement.png}{0.75\textwidth}
  \caption{Restitution de la création du dossier dans l'historique}
  \label{fig:creation-dossier}
\end{figure}


\section{Défis techniques rencontrés}

La mise en œuvre de cette logique métier, pourtant relativement simple à
énoncer module par module, s'est heurtée en pratique à plusieurs
difficultés techniques structurantes, détaillées ci-après.

\subsection{Modèles non rattachés au système de snapshot}

Certains modèles — documents, rapports, comptes bancaires — ne sont pas
directement rattachés au système de snapshot décrit précédemment. Pour
ces entités, il a fallu développer une logique de comparaison
indépendante, s'appuyant directement sur les tables d'historique
générées par \texttt{simple\_history}. Cette logique parcourt
chronologiquement chaque enregistrement d'historique afin de détecter les
transitions d'état pertinentes : création, suppression, ou changement de
visibilité d'un document.

\subsection{Relations Many-to-Many}

L'historisation des relations Many-to-Many s'est révélée
particulièrement délicate, notamment pour les labels associés aux
documents. Le mécanisme d'historisation automatique des relations M2M
fourni par \texttt{simple\_history} s'est avéré insuffisant pour ce cas
d'usage : il ne permettait pas de restituer de façon fiable et lisible
les ajouts et retraits de labels au fil du temps.

La solution retenue a consisté à créer un modèle intermédiaire explicite,
\texttt{DocumentLabel}, disposant de son propre historique. Cette
approche, bien que plus verbeuse à mettre en place, offre un contrôle
total sur la manière dont les associations sont historisées et restituées.

\subsection{Résolution des données de référence}

Les changements bruts capturés par \texttt{django-simple-history}
contiennent des identifiants techniques (clés étrangères) plutôt que des
libellés métier directement compréhensibles par un utilisateur final. Un
changement de secteur d'activité, par exemple, n'apparaît dans
l'historique brut que sous la forme d'un identifiant numérique.

J'ai donc développé une logique de résolution systématique, transformant
chaque identifiant en son libellé métier correspondant (secteurs
d'activité, codes NAF, codes ROME, catégories socio-professionnelles,
pays, etc.) avant transmission des données au frontend, garantissant un
affichage directement exploitable par les utilisateurs.

\subsection{Cohérence entre données dupliquées}

Au cours du développement, j'ai identifié un bug affectant certains
champs relatifs aux employés (fonction, catégorie socio-professionnelle,
niveau de formation) : ces champs étaient accidentellement réinitialisés
à \texttt{NULL} lors de mises à jour partielles ne les concernant pas. La
cause provenait d'une incohérence entre la structure du payload envoyé
par le frontend et celle attendue côté backend. Ce bug, bien que sans
lien direct avec l'historisation, a été détecté et corrigé grâce à
l'analyse des tables d'historique, qui révélaient clairement les
réinitialisations non désirées.

\subsection{Fusion d'événements liés}

Une même action utilisateur peut déclencher plusieurs sauvegardes
techniques distinctes côté backend (par exemple lors de la mise à jour
en cascade de plusieurs entités liées). Sans traitement particulier,
l'historique présenté à l'utilisateur se serait retrouvé fragmenté en de
multiples entrées correspondant en réalité à une seule action métier.

J'ai donc mis en place une logique de regroupement des événements
survenus dans une fenêtre de temps rapprochée, afin de présenter une vue
cohérente et lisible de l'historique, alignée avec la perception qu'a
l'utilisateur de ses propres actions.

\section{Développement frontend}

\subsection{Composant réutilisable \texttt{tab-changes}}

Côté frontend, j'ai conçu un composant Angular réutilisable,
\texttt{tab-changes}, chargé d'afficher l'historique des modifications
sous forme de timeline chronologique, à l'aide du composant PrimeNG
\texttt{p-timeline}.

\begin{figure}[H]
  \centering
  \imgw{screenshot_historique.png}{0.8\textwidth}
  \caption{Aperçu du composant d'historique final}
  \label{fig:historique}
\end{figure}

\subsection{Regroupement thématique et labels}

Les modifications sont regroupées par thématique (formation, montage
financier, stagiaires, tiers bénéficiaires, règlements, documents), avec
un système de correspondance associant chaque champ technique à son
libellé métier affiché à l'utilisateur.

\subsection{Distinction visuelle des types de modification}

Un code couleur permet de distinguer immédiatement les ajouts, les
modifications et les suppressions, facilitant la lecture rapide de
l'historique par les utilisateurs métier.

\subsection{Gestion des cas particuliers d'affichage}

Une attention particulière a été portée à la gestion de cas d'affichage
spécifiques : valeurs numériques, dates, booléens, ainsi que des objets
composites tels qu'un RIB accompagné de son justificatif, ou un organisme
de formation disposant de plusieurs attributs à afficher conjointement.

\subsection{Réutilisabilité de l'architecture}

Le pattern d'historique ainsi développé a ensuite été répliqué sur
plusieurs entités du système : demandes de financement, structures
adhérentes, comptes collaborateurs et employés. Cette réplication a
démontré la solidité et la réutilisabilité de l'architecture mise en
place, tant côté backend (\texttt{HistoryService}, snapshots) que côté
frontend (composant \texttt{tab-changes}).

\section{Méthodologie de débogage}

Face aux anomalies rencontrées durant le développement, j'ai adopté une
démarche de diagnostic systématique, consistant à vérifier l'ensemble de
la chaîne technique afin d'isoler précisément la source d'un
dysfonctionnement :

\begin{enumerate}
  \item vérification des données au niveau de la \textbf{base de
        données} elle-même, à l'aide de DBeaver ;
  \item inspection du \textbf{service backend} concerné
        (\texttt{HistoryService} ou logique associée) ;
  \item contrôle du \textbf{serializer DRF} responsable de la mise en
        forme des données exposées par l'API ;
  \item validation du contenu exact des échanges via l'\textbf{onglet
        réseau} des outils de développement du navigateur ;
  \item inspection du \textbf{composant frontend} recevant et affichant
        les données.
\end{enumerate}

Le \textbf{shell Django interactif} a également constitué un outil
précieux, me permettant d'inspecter directement l'état des données et
des tables d'historique en base, sans passer par l'interface graphique.
